% =============================================================================
% Appendix B — Proofs and Derivations
% =============================================================================

\chapter{Proofs and Derivations}
\label{app:proofs}

\section{Derivation of the OOM Transition Operator}
\index{OOM!derivation}

We derive the OOM transition operator $A^{(o,a,o',a')}$ from the
QHMM dynamics. Starting from the filtered density matrix
$\tilde\rho_t = \mathcal{E}\bigl(\Phi^{(a_t)}_{o_t}(\rho_{t-1})\bigr)$,
the HS vector evolves as:
\index{Hilbert-Schmidt!vectorization}
\index{Kraus operators!composition}

\begin{align}
v_{t+1}[\mu]
  &= \langle \mathcal{P}_\mu, \tilde\rho_{t+1} \rangle_{\mathrm{HS}} \\
  &= \Bigl\langle \mathcal{P}_\mu,
      \sum_{o'} \Phi^{(a_{t+1})}_{o'} \!\bigl(
        \mathcal{E}(\Phi^{(a_t)}_{o_t}(\mathcal{P}_\nu)) \bigr)
    \Bigr\rangle_{\mathrm{HS}} \\
  &= \sum_{o'} \operatorname{tr}\!\bigl[
      \mathcal{P}_\mu^\dagger \,
      \Phi^{(a_{t+1})}_{o'} \!\bigl(
        \mathcal{E}(\Phi^{(a_t)}_{o_t}(\mathcal{P}_\nu)) \bigr)
    \bigr] \\
  &= \sum_{o'} A^{(o_t,a_t,o',a_{t+1})}_{\mu,\nu} \, v_t[\nu],
\end{align}
which recovers Eq.\ \eqref{eq:oom_update} of Chapter~\ref{chap:background}.

\index{OOM!derivation}

\section{TD($\lambda$) Convergence}
\index{TD($\lambda$)!convergence}

The combined trace update \eqref{eq:a_update} is a stochastic gradient
descent step on the TD error objective. Under the standard assumptions of
stochastic approximation \citep{sutton2000review}—specifically bounded
rewards, a contracting value function, and a step-size schedule satisfying
$\sum \eta_t = \infty$, $\sum \eta_t^2 < \infty$—the update converges
almost surely to a fixed point of the TD($\lambda$) operator. The
combining of forward and backward traces does not affect the fixed point
(which is determined by the Bellman operator) but does affect the
variance of the learning trajectory.

\index{Bellman operator}
\index{stochastic approximation}

\section{Choi Matrix Trace-Preservation}
\index{Choi matrix!trace preservation}

For a channel $\Phi$ with Kraus operators $\{K_l\}$, the Choi matrix is
\index{trace preservation!Choi matrix}
\begin{equation}
J(\Phi) = \sum_l \ket{K_l}\!\bra{K_l}.
\end{equation}
The partial trace over the output system is
\index{partial trace}
\begin{align}
\operatorname{tr}_{\mathrm{out}}[J(\Phi)]
  &= \sum_{l} \operatorname{tr}_{\mathrm{out}}\bigl[\ket{K_l}\!\bra{K_l}\bigr] \\
  &= \sum_l (\operatorname{tr}_{\mathrm{out}} \ket{K_l}\!\bra{K_l}) \\
  &= \sum_l K_l^\dagger K_l \\
  &= I_{\mathrm{in}},
\end{align}
where the last line uses the Kraus normalization condition
$\sum_l K_l^\dagger K_l = I$. This establishes that the constraint
\eqref{eq:cptp_constraints}(a) is necessary and sufficient for
trace-preservation.

\index{Kraus normalization}
\index{CPTP!necessary and sufficient}

\section{Numerical Stability of the Forward-Backward Algorithm}
\index{forward-backward!numerical stability}

The forward recurrence $v_{t+1} = A^{(o_t,a_t,o_{t+1},a_{t+1})} v_t$ can
produce numerically vanishing or exploding values for long trajectories.
To mitigate this, we normalize at each step:
\index{log-sum-exp trick}
\index{numerical stability}
\begin{equation}
\log Z_t = \log\!\bigl(\mathbf{1}^\top v_t\bigr),
\qquad
\alpha_t = \frac{v_t}{\mathbf{1}^\top v_t},
\qquad
\tilde v_{t+1} = A^{(o_t,a_t,o_{t+1},a_{t+1})} \alpha_t,
\end{equation}
which is implemented in the code using the log-sum-exp trick to avoid
floating-point overflow. The backward pass uses the same normalization
at each step, ensuring that $\beta_t$ remains in a numerically stable
range.

\index{log-sum-exp}
\index{forward-backward!scaling}
